% ============================================================
% OP-Q11 CONVERSE CALCULATION — momentum-space OS test
% Version: v1 (2026-06-10). STATUS: DRAFT FOR GPT VERIFICATION.
% Companion to OP_Q11_proposal_v2.tex. NOT inserted into preprint.
% ============================================================

\section*{Momentum-space OS test for the free minimal Euclidean
Dirac kernel}

% ------------------------------------------------------------
% 0. Euclidean conventions (NEW — not yet fixed in main text)
% ------------------------------------------------------------
\subsection*{0. Conventions}
Euclidean coordinates $x=(w,\vec x)$ on $\mathbb R^4$; reflection
$\Theta:(w,\vec x)\mapsto(-w,\vec x)$, consistent with the
ordered-branch convention $t=w$.
Hermitian Euclidean gamma matrices
\[
\{\gamma_A,\gamma_B\}=2\delta_{AB},\qquad
\gamma_A^{\dagger}=\gamma_A,\qquad A,B\in\{0,1,2,3\},
\]
with $\gamma_0$ associated to the $w$-direction.
[NOTE: the main text (\S sec:dirac\_eq) currently fixes only the
Lorentzian convention $\{\gamma^\mu,\gamma^\nu\}=2\eta^{\mu\nu}$;
the Euclidean block above is proposed for adoption and must be
reconciled at implementation time.]
Minimal first-order kernel: Euclidean Dirac operator and two-point
kernel
\[
D_E=\gamma_A\partial_A+m,\qquad m>0,\qquad
S_E=D_E^{-1},\qquad
\tilde S_E(p)=\frac{-i\gamma_A p_A+m}{p^2+m^2}.
\]
Spinorial reflection matrix: $\gamma_0$, acting as
$(\Theta\psi)(w,\vec x)=\gamma_0\,\psi(-w,\vec x)$, anti-linearly
on coefficients; for Grassmann polynomials $\Theta$ additionally
reverses factor order (graded OS form).
Notation: $\omega(\vec p)=\sqrt{\vec p^{\,2}+m^2}\ge m$.
Throughout, the ECT normalisation $e^{-S_E/S_0}$ only rescales the
kernel by a positive constant and is suppressed.

% ------------------------------------------------------------
% 1. Half-space contour evaluation (common to both directions)
% ------------------------------------------------------------
\subsection*{1. Reflected kernel at fixed spatial momentum}
Take a spinor test function $g(w,\vec x)$ supported in $w>0$ and
define the OS two-point form
\[
Q[g]\;=\;\int d^4x\,d^4y\;
\overline{(\Theta g)(x)}\;K(x-y)\;g(y),
\]
where $K$ is the relevant two-point kernel (specified per
direction below) and the bar includes the spinorial reflection
matrix.
Writing the reflection explicitly, the $w$-dependence enters only
through $w_x+w_y>0$.
The $p_0$-integral is evaluated by closing the contour in the lower
half-plane (factor $e^{-ip_0(w_x+w_y)}$, $w_x+w_y>0$), picking the
pole $p_0=-i\omega$:
\[
\int\!\frac{dp_0}{2\pi}\,
\frac{e^{-ip_0(w_x+w_y)}\,(-i\gamma_0 p_0-i\vec\gamma\!\cdot\!\vec p+m)}
{p_0^2+\omega^2}
=\frac{e^{-\omega(w_x+w_y)}}{2\omega}\,
\bigl(m-\omega\gamma_0-i\vec\gamma\!\cdot\!\vec p\bigr).
\]
Hence, at fixed $\vec p$,
\[
Q[g]=\int\!\frac{d^3p}{(2\pi)^3}\;
h(\vec p)^{\dagger}\,\mathcal M(\vec p)\,h(\vec p),
\qquad
h(\vec p)=\int_0^{\infty}\!dw\,e^{-\omega w}\,
\tilde g(w,\vec p),
\]
up to positive normalisation, where $\mathcal M(\vec p)$ is the
spinor-space matrix obtained after attaching the reflection matrix
$\gamma_0$:
\[
\mathcal M(\vec p)\;=\;
\frac{1}{2\omega}\,\gamma_0
\bigl(m-\omega\gamma_0-i\vec\gamma\!\cdot\!\vec p\bigr)\,
[\,\text{ordering factor}\,].
\]
The ordering factor is where the statistics enters: it is $(-1)$
from the Grassmann order reversal in the graded form, and $(+1)$
for commuting fields.
[Sign bookkeeping to be verified jointly with the adjoint
convention; the relative sign between the two quantisations is the
invariant content.]

% ------------------------------------------------------------
% 2. Direct direction: Grassmann construction is RP
% ------------------------------------------------------------
\subsection*{2. Grassmann case: projector identity}
With the Grassmann ordering factor, the fixed-$\vec p$ matrix is
\[
\mathcal M_{\rm Gr}(\vec p)
=\frac{1}{2\omega}\,
\bigl(\omega+ m\gamma_0+i\gamma_0\vec\gamma\!\cdot\!\vec p\bigr).
\]
Define $A=m\gamma_0+i\gamma_0\vec\gamma\!\cdot\!\vec p$.
Using $\{\gamma_0,\vec\gamma\}=0$ and hermiticity:
$A^{\dagger}=A$ and
\[
A^2=m^2+\vec p^{\,2}=\omega^2
\quad\Longrightarrow\quad
\mathcal M_{\rm Gr}=\frac{\omega+A}{2\omega},
\qquad
\mathcal M_{\rm Gr}^2=\mathcal M_{\rm Gr},
\qquad
\mathcal M_{\rm Gr}^{\dagger}=\mathcal M_{\rm Gr}.
\]
$\mathcal M_{\rm Gr}(\vec p)$ is an orthogonal projector (rank~2:
$\mathrm{tr}\,\mathcal M_{\rm Gr}=2$), hence
$\mathcal M_{\rm Gr}\ge 0$ and $Q[g]\ge 0$ for all admissible $g$:
the graded OS form is reflection positive.
[This reproduces the standard OS fermionic construction in ECT
conventions; it is the spinorial analogue of the
$e^{-\omega(w_x+w_y)}$ rank-one positivity used for the scalar
subsectors in \S sec:qs\_reflection\_positivity.]

% ------------------------------------------------------------
% 3. Converse: commuting/symmetric quantisation fails RP
% ------------------------------------------------------------
\subsection*{3. Commuting case: explicit negative modes}
For commuting fields the two-point Schwinger function must be
permutation-symmetric; for the \emph{same} kernel the reflected
fixed-$\vec p$ form is therefore built from the symmetrised kernel.
Symmetrisation removes the part of
$\gamma_0(m-\omega\gamma_0-i\vec\gamma\!\cdot\!\vec p)$ that is odd
under the combined exchange, which is precisely the
$i\gamma_0\vec\gamma\!\cdot\!\vec p$ term (odd in $\vec p$ together
with the spinor transposition), leaving
\[
\mathcal M_{\rm sym}(\vec p)
=\frac{1}{2\omega}\,\bigl(m\gamma_0-\omega\bigr).
\]
Diagonalisation at fixed $\vec p$: $\gamma_0$ has eigenvalues
$\pm1$ (each twice), so
\[
\mathrm{spec}\,\mathcal M_{\rm sym}(\vec p)
=\Bigl\{\frac{m-\omega}{2\omega},\;
\frac{-m-\omega}{2\omega}\Bigr\},
\]
each with multiplicity two.
Since $\omega\ge m$ with equality only at $\vec p=0$:
\[
\frac{-m-\omega}{2\omega}\le-\frac{m}{\omega+m}<0
\quad\text{for all }\vec p,
\qquad
\frac{m-\omega}{2\omega}<0
\quad\text{for }\vec p\neq0.
\]
\emph{Explicit negative-norm test function.}
Choose any $\vec p_\ast$ (including $\vec p_\ast=0$), a spinor
$u_-$ with $\gamma_0 u_-=-u_-$, and
$g(w,\vec x)=\chi(w)\,e^{i\vec p_\ast\cdot\vec x}
\rho(\vec x)\,u_-$ with $\chi$ smooth, supported in $w>0$, and
$\rho$ a broad spatial envelope concentrating $\tilde g$ near
$\vec p_\ast$.
Then
\[
Q[g]\;\longrightarrow\;
\frac{-m-\omega(\vec p_\ast)}{2\omega(\vec p_\ast)}\;
\bigl|h(\vec p_\ast)\bigr|^2\;<\;0 .
\]
The OS form is not merely indefinite but strictly negative on the
$\gamma_0=-1$ eigenmodes; reflection positivity fails for the
commuting/symmetric quantisation of the same minimal kernel.

\emph{Degenerate alternative (second horn).}
If one instead insists on a genuine commuting Gaussian measure, the
covariance must itself be symmetric, which deletes the odd
first-order part $-i\gamma\!\cdot\!p$ entirely and reduces the
kernel to $m/(p^2+m^2)$: a scalar-type kernel in spinor clothing,
no longer the minimal first-order Dirac kernel.
Either horn —--- keep the kernel and lose positivity, or keep
symmetry and lose the kernel --- excludes the commuting
quantisation of the half-integer minimal kernel.

% ------------------------------------------------------------
% 4. Caveats / verification list for GPT
% ------------------------------------------------------------
\subsection*{4. Caveats and verification list}
\begin{enumerate}[label=(\alph*)]
  \item Sign/ordering bookkeeping in Section~1 (the ``ordering
    factor'' and the anti-linear adjoint) is the most
    error-prone step; the invariant claim is the \emph{relative}
    sign between the Grassmann and commuting forms.
  \item The symmetrisation step in Section~3 (which part of the
    kernel is exchange-odd) must be re-derived with explicit
    spinor-transpose conventions (charge-conjugation matrix), not
    only by the $\vec p$-parity shortcut used here.
  \item Conventions are proposed, not yet adopted: the main text
    has no Euclidean gamma block (\S sec:dirac\_eq is Lorentzian).
  \item The computation covers the minimal first-order kernel only
    (scope restriction of Q11.1a).
  \item Massless limit $m\to0$: the $\vec p=0$ symmetric eigenvalue
    degenerates to $0$; the strict-negativity statement should be
    quoted for $m>0$ (mass gap assumption of Q11.1a).
\end{enumerate}
% ============================================================
